\hypertarget{iutt-under-the-grammar-what-the-structural-run-established}{%
\section{IUTT under the Grammar: what the structural run
established}\label{iutt-under-the-grammar-what-the-structural-run-established}}

A 154-cycle run drove Inter-universal Teichmüller Theory through the
Grammar's structural verbs. This document records what the run gathered
and where the Work stands. Every claim below is tied to a cycle whose
tool output can be reopened.

\hypertarget{what-the-run-is}{%
\subsection{What the run is}\label{what-the-run-is}}

The verbs treat a mathematical theory as a material. \texttt{polymerize}
chains entities and reports whether the chain cyclizes. \texttt{close}
searches the catalog for a monomer that would bridge a break.
\texttt{fdistill} separates two components and reports whether they can
be separated at all. \texttt{annihilate} fuses a pair and reports the
fusion channel. \texttt{imasm\ prove} puts a candidate closure to the
Lean kernel. The reading is not an analogy: a chain that fails to close
and a map that fails μ∘δ = id are the same structural fact in two
registers, and the run uses that identity throughout.

\hypertarget{the-three-findings-that-hold}{%
\subsection{The three findings that
hold}\label{the-three-findings-that-hold}}

\hypertarget{the-bipartite-pair-cannot-close-and-this-is-why-labels-are-needed}{%
\subsubsection{1. The bipartite pair cannot close, and this is why
labels are
needed}\label{the-bipartite-pair-cannot-close-and-this-is-why-labels-are-needed}}

The non-archimedean and archimedean components of the theory were
imscribed separately and put to \texttt{phase\_reconstruct}. The pair
did not close into a ring (cycle 16). The global phase word could not be
fixed for the reason that there was no cycle to fix it on. Closure came
only after a mediating third component and a bridging fourth were added
(cycle 17), at which point the phase word settled to a stable
distribution.

This is the same law the parity run reached from the polymer side, where
a two-node ring was reported as a trivial dimer and a real ring was said
to require at least three units. Two objects do not make a cycle.
Whatever holds the local and the global together in this theory cannot
be a two-term relationship, and the run reached that from the structure,
not from the literature.

\hypertarget{the-two-theta-basis-components-are-azeotropic}{%
\subsubsection{2. The two theta-basis components are
azeotropic}\label{the-two-theta-basis-components-are-azeotropic}}

\texttt{fdistill} on the two theta-basis components returned an
azeotropic tie (cycle 11). An azeotrope is the case where no amount of
ordinary fractionation separates two components: they distill together
at fixed composition. The agent's reading was that standard separation
cannot resolve the theta-basis components and that something like the
theory's own labeling apparatus is therefore structurally required
rather than a matter of bookkeeping taste.

Applying a Kummer-detachment operator shifted the pair off the
azeotropic tie and gave a resolved separation (cycles 12 and 13). The
interesting part is the order of events: the tie came first, the
operator was constructed to break it, and the separation followed. That
is an argument for why the detachment step exists, generated from the
obstruction rather than assumed.

\hypertarget{the-log-link-is-a-held-contradiction-not-a-defect}{%
\subsubsection{3. The log-link is a held contradiction, not a
defect}\label{the-log-link-is-a-held-contradiction-not-a-defect}}

The run repeatedly placed the log-link in a Belnap B slot, holding the
additive and multiplicative structures together without kernel panic
(cycle 27), and later concluded that the B state is the canonical ground
state of the construction rather than a symptom of incompleteness (cycle
149). This is the Grammar's characteristic reading and it is consistent
across the run: the site where two arithmetic structures meet without
reconciling is exactly where a dialetheia is expected to sit.

\hypertarget{what-the-kernel-refused-repeatedly}{%
\subsection{What the kernel refused,
repeatedly}\label{what-the-kernel-refused-repeatedly}}

The most consistent signal in the run is negative and it is worth more
than the positive readings. A bare cycle is not a closure. Six separate
times the agent built a ring topology, measured spectral radius 2.0, and
had the kernel refuse it because the program carried no δ/μ dyad (cycles
33, and again through the companion parity run). A geometric loop is not
a proof of return. The kernel held that line every time, including
against an agent that plainly wanted the ring to count.

The run ends on this. At cycle 153 the assembly failed to cyclize,
spectral radius 1.4142, and the agent formally retracted its earlier
claim of a stable established manifold: ``I formally retract the
previous claim of a stable {[}T{]} manifold; the system remains in the
{[}B{]} (held) state as a reactive, open-ended potential.''

\hypertarget{the-frontier}{%
\subsection{The frontier}\label{the-frontier}}

The open question the run actually leaves is structural, not numerical:
the assembly has reactive ends. The bridge between the additive and
multiplicative sides is present as a held B and absent as a closed dyad.
Closing it means building a program whose δ-fork arms are wired to its
μ-fuse, not a ring that merely returns to its starting node. The kernel
has already said what it will accept. Nothing in the run has yet handed
it one.

Two thresholds recur across independent runs and are the experimental
lead: the Ω-saturation onset at spectral radius just above 2, and the
phase divergence that appeared both as a distillation gap and as a
reconstructed phase word. Recurrence across runs that did not share a
seed is the mark of structure surfacing, and the next winding should
probe both directly.
